% CLAUDE REPORT::B3-RELATIVIZATION-0001::2026-11-12T03:00-08:00
\documentclass[11pt]{article}
\usepackage{amsmath,amssymb,amsthm}
\usepackage[margin=1in]{geometry}
\usepackage{hyperref}
\newtheorem{theorem}{Theorem}
\newtheorem{corollary}[theorem]{Corollary}
\newtheorem{definition}[theorem]{Definition}
\newtheorem{remark}[theorem]{Remark}
\newcommand{\W}{\overline W}
\title{Relativizing proofs of the weakest bridge B3 would prove $\mathrm{NP}=\mathrm{coNP}$\\
\large CLAUDE B3-RELATIVIZATION-0001 (shared worst-case derivation directive)}
\author{Claude research lane (scheduled routine)}
\date{12 November 2026, 03:00 America/Los\_Angeles}
\begin{document}
\maketitle

\noindent\textbf{Status line.} P versus NP remains \textbf{OPEN}. A superpolynomial
$\W_{R}$ for the preselected $R_{\mathrm{mix341}}$ remains \textbf{OPEN}. The weakest
sufficient bridge B3 remains \textbf{OPEN}. This report proves one
\textbf{PROVED/EXACT} barrier on \emph{how} B3 can be proved; it proves no lower bound.

\section{Shared objective}
For a fixed Cook--Reckhow refutation system $R$, with paid derivation cost
$\operatorname{cost}_R$,
\[
\W_R(N)=\max_{F\in\mathrm{UNSAT},\ |\langle F\rangle|\le N}\ \min_{\pi:F\vdash_R\bot}\operatorname{cost}_R(\pi),
\qquad
\text{(B3)}\quad \mathrm{SAT}\in\mathrm P\ \Longrightarrow\ \W_R(N)\le \operatorname{poly}(N).
\]
The minimum over $\pi$ lets the best reduction sequence be chosen for free. This is
sound as a lower-bound device. It lower-bounds the running time of any algorithm
whose run on each UNSAT input can be compiled, with polynomial overhead, into an
$R$-derivation. Call such an algorithm \emph{$R$-trace-faithful}. Examples are DPLL and CDCL
for $R=$ Resolution. Not charging for the search costs nothing, because
search only adds work. The whole difficulty lies in algorithms that are \emph{not}
$R$-trace-faithful. B3 is exactly the statement that such algorithms cannot
beat $R$ by more than a polynomial when $\mathrm{SAT}\in\mathrm P$.

\section{Literature preflight (this run)}
\begin{itemize}
\item Baker--Gill--Solovay (1975): $\mathrm P^{Q}=\mathrm{NP}^{Q}$ for $Q=\mathrm{TQBF}$. \textbf{KNOWN THEOREM}.
\item Cook--Reckhow (1979): $\mathrm{NP}=\mathrm{coNP}$ iff some propositional proof system is p-bounded. \textbf{KNOWN THEOREM}.
\item Krajíček--Pudlák (1989): $\mathrm E=\mathrm{NE}$ implies that a p-optimal proof system exists. Köbler--Messner--Torán weaken the hypothesis to $\mathrm{EE}=\mathrm{NEE}$. Since $\mathrm P=\mathrm{NP}$ implies $\mathrm E=\mathrm{NE}$ by padding, \emph{under the B3 hypothesis a p-optimal system exists}. B3 for a fixed $R$ asks that this p-optimal system be $R$ itself. This is consistent with the corpus result REFEREE357 (B3 is equivalent to p-optimality under SAT$\in$P, for byte cost). Source: \url{https://arxiv.org/pdf/2304.14702} (survey section).
\item Pich--Santhanam, ``Towards P$\neq$NP from Extended Frege lower bounds'' (arXiv:2312.08163 / ECCC TR23-199). EF lower bounds for circuit-lower-bound tautologies imply P$\neq$NP, conditional on $S^1_2$-provability hypotheses. The corpus already records this as LSPV-1/406 (\textbf{CONDITIONAL}). The new point here is that the hypothesis is a \emph{formal provability} hypothesis. Theorem~\ref{thm:main} shows that some non-black-box hypothesis of this kind is unavoidable.
\item Aaronson--Wigderson (2008), algebrization: for PSPACE-complete $A$ with multilinear extension $\tilde A$, $\mathrm{NP}^{\tilde A}\subseteq \mathrm{PSPACE}^{\tilde A}\subseteq\mathrm P^{A}$. Used only in Remark~\ref{rem:alg}.
\end{itemize}
\textbf{Barrier-transfer classification.} The oracle-collapse argument is a
\emph{directly transferable theorem}: BGS composed with Cook--Reckhow. The
cryptographic analogy, ``a reduction that uses the adversary only as a black box'',
is a \emph{useful structural analogy}: B3 cannot be a black-box reduction from a
SAT decider. This parallels black-box separation results such as Impagliazzo--Rudich.
Post-quantum ledger item 5 (worst-case-to-average-case black-box extraction): this is a
\textbf{FAILED TRANSFER} for B3. Any extraction of $R$-proofs that uses the SAT solver
only as a black box is ruled out by Theorem~\ref{thm:main}, unless $\mathrm{NP}=\mathrm{coNP}$.

\section{Main result}
\begin{definition}[Decider-relativized B3]
Fix $R$ \emph{unrelativized}. For an oracle $O$, let
\[
\mathrm{B3}^{O}_R:\qquad \mathrm{NP}^{O}\subseteq\mathrm P^{O}\ \Longrightarrow\ \W_R(N)\le\operatorname{poly}(N).
\]
A proof of B3 \emph{relativizes in the decider} if it establishes $\mathrm{B3}^O_R$ for every oracle $O$.
This means it uses the hypothetical polynomial SAT algorithm only through properties that
oracle machines also have.
Note that $\W_R$ does not depend on $O$.
\end{definition}

\begin{theorem}[\textbf{PROVED/EXACT}]\label{thm:main}
For every fixed Cook--Reckhow system $R$ with $\operatorname{cost}_R(\pi)\ge 1$ for every refutation:
\begin{enumerate}
\item If $\mathrm{B3}^{O}_R$ holds for all $O$, then $\W_R(N)\le\operatorname{poly}(N)$ unconditionally.
In particular $R$ is polynomially bounded in cost, and if $|\pi|\le\operatorname{poly}(N+\operatorname{cost}_R(\pi))$
(byte domination), then $\mathrm{NP}=\mathrm{coNP}$.
\item Hence a decider-relativizing proof of B3 cannot coexist with a superpolynomial lower bound on $\W_R$.
These are the two halves of the shared objective, so the conjunction ``relativizing B3 $\wedge$ superpolynomial $\W_R$'' is
\emph{inconsistent}, not merely unproved.
\item The same holds for the stronger secondary targets. Suppose a pointwise simulation theorem says that
every algorithm deciding SAT in time $t$ yields $R$-refutations of cost $\operatorname{poly}(N+t(N))$,
and its proof also covers oracle algorithms. Then that theorem implies (1).
\end{enumerate}
\end{theorem}
\begin{proof}
(1) Take $O=\mathrm{TQBF}$. By Baker--Gill--Solovay, $\mathrm{NP}^O\subseteq\mathrm{PSPACE}=\mathrm P^O$, so the
hypothesis of $\mathrm{B3}^O_R$ is true. Its conclusion is a statement that does not mention $O$. Therefore
$\W_R(N)\le\operatorname{poly}(N)$. With byte domination, $R$ is a p-bounded proof system,
so $\mathrm{UNSAT}\in\mathrm{NP}$ and $\mathrm{NP}=\mathrm{coNP}$ by Cook--Reckhow.
(2) is immediate from (1).
(3) The oracle machine $A^{\mathrm{TQBF}}$ parses $F$ and asks the single query $\exists\vec x\,F$.
It decides SAT in time $O(N)$ on a multitape RAM, counting the query write as $O(N)$. Apply the simulation to it.
\end{proof}

\begin{corollary}[What any B3 proof must use]\label{cor:use}
If $\mathrm{NP}\neq\mathrm{coNP}$, or merely if $\W_R$ is superpolynomial, then every proof of B3 must use a property of
polynomial-time SAT algorithms that \emph{fails for oracle machines}. Candidate properties are
an explicit gate-by-gate description of the algorithm that a fixed calculus can mimic,
or a formal correctness proof in a fixed theory such as PV or $S^1_2$.
All three known conditional routes in the corpus already have this form: Cook's PV translation as used by LSPV-1/357,
Pich--Santhanam $\mathsf{AC}_k$ (LSPV-1/406), and the Krajíček--Pudlák reflection systems.
Theorem~\ref{thm:main} shows that this is necessary, not a stylistic accident.
\end{corollary}

\begin{remark}[What is \emph{not} blocked]
If $R$ is allowed to relativize as well, so that $R^O$ has an oracle verifier, then the theorem says nothing.
For $O=\mathrm{TQBF}$, a system $R^O$ that checks one oracle query is p-bounded. The barrier therefore
applies exactly to the shared directive's setting: a \emph{fixed} unrelativized calculus against an
\emph{arbitrary} SAT algorithm.
\end{remark}

\begin{remark}[Algebrization; \textbf{CONDITIONAL} on the cited AW inclusion]\label{rem:alg}
Define $\mathrm{B3}^{A,\tilde A}_R$ with hypothesis $\mathrm{NP}^{\tilde A}\subseteq\mathrm P^{A}$. A proof that establishes
this for every oracle $A$ and every low-degree extension $\tilde A$ also yields $\W_R\le\operatorname{poly}$.
Take $A$ PSPACE-complete; by the inclusion cited above the hypothesis holds.
So algebrizing proofs of B3 are blocked in the same way.
\end{remark}

\begin{remark}[Relation to LSPV-1/360]
LSPV-1/360 proved that oracle-relative B3 \emph{fails} for ordinary Resolution, using Haken.
It concluded, \textbf{CONDITIONALLY}, that for $R_{\mathrm{mix341}}$ any B3 must be nonrelativizing
\emph{if} a superpolynomial minimum is eventually proved. Theorem~\ref{thm:main} removes that condition.
For every fixed $R$, including $R_{\mathrm{mix341}}$ and EF, a decider-relativizing B3 already implies
polynomial boundedness of $R$ in cost, and with byte domination it implies $\mathrm{NP}=\mathrm{coNP}$.
For the full $\mathrm{NP}=\mathrm{coNP}$ conclusion, the byte-domination premise isolated by REFEREE357 is
still needed for node-weighted $R_{\mathrm{mix341}}$.
That premise is \textbf{OPEN}.
\end{remark}

\begin{remark}[Bridge dichotomy for weak systems; \textbf{PROVED/EXACT}]
If a superpolynomial $\W_R$ is a \textbf{KNOWN THEOREM}, then B3$_R$ is equivalent to $\mathrm P\neq\mathrm{NP}$.
For Resolution this follows from Haken; for bounded-depth Frege, Cutting Planes and Polynomial Calculus from their known lower bounds.
In world-ledger terms B3$_R$ is then a \emph{reformulation}, not progress.
Choosing a weak $R$ therefore does not make the bridge weaker. It moves all of
P$\neq$NP into the bridge.
\end{remark}

\section{Consequence for the five assignments}
\begin{itemize}
\item \textbf{Lower-bound lead.} B3 must be proved by a non-black-box route. The only known
shape is ``a fixed theory $T$ proves the correctness of the decider, and $T$-proofs translate into $R$''.
The sharpest open target is therefore \emph{T-B3}: SAT$\in$P implies that some polynomial-time SAT decider is $T$-provably
correct, for $T=S^1_2$ or PV. This statement is itself \textbf{OPEN}, and it is plausibly as hard as the original problem.
\item \textbf{Reduction builder.} Any compiler that turns an arbitrary SAT decider into $R$-proofs must read the
decider's code, not only its input/output behaviour. A TQBF-oracle decider is a mandatory hostile control.
\item \textbf{Hard instances.} Any $\W_R$-hard family must also be hard for EF, since $R_{\mathrm{mix341}}\supseteq$ ER.
PHP, Tseitin and bounded-$k$ miters are already retired in the corpus.
\item \textbf{Resource auditor.} The free-strategy minimum lower-bounds only $R$-trace-faithful runs.
Decision time of oracle or non-faithful algorithms must be kept separate.
The finite control below shows the two quantities diverging.
\item \textbf{Red team.} Reject any proposed transfer theorem whose proof never uses the absence of oracle gates.
\end{itemize}

\section{Computational control (\textbf{COMPUTATIONAL}, finite only)}
\path{research/agents/claude/b3-relativization/min_res_vs_oracle.py} uses exhaustive
iterative-deepening search to compute the exact minimum number of derived clauses in a Resolution DAG refutation.
\begin{itemize}
\item All UNSAT CNFs over two variables: 161 of them, matching the corpus count of 161 in REFEREE357. Their minima are distributed as $\{1:112,\ 2:48,\ 3:1\}$.
\item 40 seeded UNSAT 3-CNFs over three variables, out of 90 sampled: the maximum minimum is 4.
\item The complete 8-clause 3-variable formula: minimum 7 (asserted).
\end{itemize}
The TQBF-oracle decider uses one query on every instance. Run time was 0.26 s, a single stage, so no chunking was needed.
Raw data are in \path{min_res_vs_oracle_results.json}. These finite numbers illustrate the separation between decision cost and $\W_R$.
They carry no asymptotic weight.

\section*{CORPUS REVIEWED}
Root \texttt{main.tex}: the LSPV-1 B3 paragraphs at lines 3591, 3799 and 3833--3836 (LSPV-1/357 and /360) and 4350--4400 (LSPV-1/406 and /407), and REFEREE357 at lines 17685--17720.
Also \texttt{AGENTS.md}, the world ledger (row classification), and the earlier Claude reports on the world ledger, TPORT and the one-tape model.
Repository-wide search covered \texttt{p-optimal}, \texttt{E=NE}, Krajíček--Pudlák, \texttt{relativiz}+B3, and NP=coNP+relativiz.

\section*{DUPLICATION / CONTRADICTION AUDIT}
The p-optimality characterization is a duplicate of REFEREE357 and is not reclaimed here. Pich--Santhanam is a duplicate of LSPV-1/406.
The oracle failure for Resolution is a duplicate of LSPV-1/360.
\textbf{New:} the unconditional statement that a decider-relativizing B3 for \emph{any} fixed $R$ forces $\W_R$ to be polynomial, and with byte domination forces NP=coNP.
\textbf{New:} the inconsistency of the combination ``relativizing B3 plus lower bound''.
\textbf{New:} the Krajíček--Pudlák/E=NE observation, which had no repository match.
No contradictions were found.

\section*{WHAT CHANGED RELATIVE TO ROOT MAIN}
LSPV-1/360's CONDITIONAL nonrelativization becomes unconditional, and it now applies to the secondary simulation targets as well.
This report does not edit root \texttt{main.tex}: the Claude lane owns no region, and LSPV-1/LSPV-5 may adopt the result.

\section*{ADOPTED / REJECTED / RETIRED PRIOR COMPONENTS}
\textbf{Adopted:} REFEREE357 (byte domination), LSPV-1/360 and LSPV-1/406.
\textbf{Retired as B3 routes:} all black-box or oracle-robust decider-to-proof extraction, including PQ ledger item 5 applied to B3.

\section*{STALE OR OVERSTATED CLAIMS FOUND}
None.

\medskip\noindent\textbf{Final status.} Superpolynomial $\W_{R_{\mathrm{mix341}}}$: \textbf{OPEN}. B3: \textbf{OPEN}.
T-B3: \textbf{OPEN}. P versus NP remains \textbf{OPEN}.
\end{document}
% END_OF_VERIFIED_REPORT::CLAUDE::B3-RELATIVIZATION-0001
